
\begin{definition}[Full-support specialization of one fusion layer]%
    \label{def:mpdo_full_support_virtual_projector_one_fusion}
    \lean{MPOTensor.BNTFusionIsometryFamily.projectorQBlock}
    \lean{MPOTensor.BNTFusionIsometryFamily.conjugatedProjectorQBlock}
    \leanok
    \uses{def:mpdo_bnt_fusion_isometry_family}
    For a full-support fusion family, let $P_\gamma$ act on the bond space of
    $M_\gamma$.  Define
    \[
        Q^{\mathrm{fs}}_{\alpha,\beta}
          =\bigoplus_\gamma
             \chi_{\alpha,\beta,\gamma}\otimes P_\gamma,
        \qquad
        \widehat Q^{\mathrm{fs}}_{\alpha,\beta}
          =U_{\alpha,\beta}^{\dagger}
             Q^{\mathrm{fs}}_{\alpha,\beta}U_{\alpha,\beta}.
    \]
    The terminal contraction is the same as in the source operator.  The
    additional condition is that the active support is the whole product bond
    space.

    \textbf{Scope restriction (full-support fusion family):} Full support is
    not an additional hypothesis of the arbitrary-chain source statement.
    This distinction is
    documented in
    \path{docs/paper-gaps/cpsv16_topological_projector_recursion.tex}.
\end{definition}

\begin{theorem}[Projection property of the full-support specialization]%
    \label{thm:mpdo_full_support_virtual_projector_one_fusion}
    \lean{MPOTensor.BNTFusionIsometryFamily.projectorQBlock_eq_unweighted}
    \lean{MPOTensor.BNTFusionIsometryFamily.projectorQBlock_isStarProjection}
    \lean{MPOTensor.BNTFusionIsometryFamily.%
        conjugatedProjectorQBlock_isOrthogonalProjection}
    \leanok
    \uses{def:mpdo_full_support_virtual_projector_one_fusion,
        thm:mpdo_pair_fusion_bnt_completeness,
        thm:star_projection_coisometric_transport,
        thm:orthogonal_projection_star_projection_channel}
    Suppose that the labelled tensors form a basis of normal tensors and the
    positive trace-power coefficients are independent of the positive chain
    length.  If each $P_\gamma$ is a self-adjoint idempotent on the bond space,
    then $Q^{\mathrm{fs}}_{\alpha,\beta}$ is a self-adjoint idempotent and
    $\widehat Q^{\mathrm{fs}}_{\alpha,\beta}$ is an orthogonal
    projection.
\end{theorem}
\begin{proof}\leanok
    Length independence gives
    $\chi_{\alpha,\beta,\gamma}=1$, hence
    $(Q^{\mathrm{fs}}_{\alpha,\beta})^2
      =Q^{\mathrm{fs}}_{\alpha,\beta}$ and
    $(Q^{\mathrm{fs}}_{\alpha,\beta})^\dagger
      =Q^{\mathrm{fs}}_{\alpha,\beta}$.
    The basis-of-normal-tensors completeness theorem gives
    $U_{\alpha,\beta}U_{\alpha,\beta}^{\dagger}=1$, so
    \[
      (\widehat Q^{\mathrm{fs}}_{\alpha,\beta})^2
        =U_{\alpha,\beta}^{\dagger}Q^{\mathrm{fs}}_{\alpha,\beta}
          (U_{\alpha,\beta}U_{\alpha,\beta}^{\dagger})
          Q^{\mathrm{fs}}_{\alpha,\beta}U_{\alpha,\beta}
        =\widehat Q^{\mathrm{fs}}_{\alpha,\beta}.
    \]
\end{proof}

\begin{theorem}[Left iterated fusion completeness]%
    \label{thm:mpdo_left_fusion_selector_completeness}
    \lean{MPOTensor.BNTFusionIsometryFamily.%
        leftFusionIsometry_mul_conjTranspose_eq_one_of_lengthIndependent_of_selectorWords}
    \leanok
    \uses{thm:mpdo_pair_fusion_selector_completeness,
        def:mpdo_left_fusion_isometry}
    Under the hypotheses of
    Theorem~\ref{thm:mpdo_pair_fusion_selector_completeness},
    \[
        U^{\mathrm L}_{\alpha,\beta,\gamma}
        (U^{\mathrm L}_{\alpha,\beta,\gamma})^{\dagger}=1.
    \]
\end{theorem}
\begin{proof}\leanok
    Pair fusion completeness gives
    \[
        \sum_\delta r_{a,b}^{\delta}D_\delta=D_aD_b.
    \]
    Applying this equality first to $(\delta,\gamma)$ and then to
    $(\alpha,\beta)$ shows that the source and target dimensions of
    $U^{\mathrm L}_{\alpha,\beta,\gamma}$ coincide.  A square isometry is
    surjective.
\end{proof}

\begin{theorem}[Right iterated fusion completeness]%
    \label{thm:mpdo_right_fusion_selector_completeness}
    \lean{MPOTensor.BNTFusionIsometryFamily.%
        rightFusionIsometry_mul_conjTranspose_eq_one_of_lengthIndependent_of_selectorWords}
    \leanok
    \uses{thm:mpdo_pair_fusion_selector_completeness,
        def:mpdo_right_fusion_isometry}
    Under the same hypotheses,
    \[
        U^{\mathrm R}_{\alpha,\beta,\gamma}
        (U^{\mathrm R}_{\alpha,\beta,\gamma})^{\dagger}=1.
    \]
\end{theorem}
\begin{proof}\leanok
    The right-associated dimension calculation is
    \[
        \sum_{\delta,\varepsilon}
          r_{\beta,\gamma}^{\delta}r_{\alpha,\delta}^{\varepsilon}
          D_\varepsilon
        =D_\alpha(D_\beta D_\gamma).
    \]
    The right iterated fusion map is therefore a square isometry.
\end{proof}

\begin{theorem}[Right adjoint inverse for the full triple-fusion comparison]%
    \label{thm:mpdo_triple_fusion_comparison_right_inverse}
    \lean{MPOTensor.BNTFusionIsometryFamily.%
        tripleFusionComparison_mul_conjTranspose_eq_one_of_lengthIndependent_of_selectorWords}
    \leanok
    \uses{lem:mpdo_triple_fusion_comparison_left_range,
        thm:mpdo_left_fusion_selector_completeness}
    Under the same hypotheses, the comparison of the full direct sums satisfies
    \[
        C_{\alpha,\beta,\gamma}C_{\alpha,\beta,\gamma}^{\dagger}=1.
    \]
\end{theorem}
\begin{proof}\leanok
    The left range identity reduces the product to
    $U^{\mathrm L}(U^{\mathrm L})^{\dagger}$, which is the identity by
    Theorem~\ref{thm:mpdo_left_fusion_selector_completeness}.
\end{proof}

\begin{theorem}[Left adjoint inverse for the full triple-fusion comparison]%
    \label{thm:mpdo_triple_fusion_comparison_left_inverse}
    \lean{MPOTensor.BNTFusionIsometryFamily.%
        conjTranspose_mul_tripleFusionComparison_eq_one_of_lengthIndependent_of_selectorWords}
    \leanok
    \uses{lem:mpdo_triple_fusion_comparison_right_range,
        thm:mpdo_right_fusion_selector_completeness}
    Under the same hypotheses,
    \[
        C_{\alpha,\beta,\gamma}^{\dagger}C_{\alpha,\beta,\gamma}=1.
    \]
    This proves invertibility on the full direct sums.  It does not yet prove
    invertibility of a fixed-final multiplicity matrix.
\end{theorem}
\begin{proof}\leanok
    Apply the right range identity and
    Theorem~\ref{thm:mpdo_right_fusion_selector_completeness}.
\end{proof}

\begin{theorem}[Vanishing between distinct final sectors]%
    \label{thm:mpdo_triple_fusion_comparison_final_sector_vanishing}
    \lean{MPOTensor.BNTFusionIsometryFamily.%
        tripleFusionComparison_finalSector_submatrix_eq_zero}
    \leanok
    \uses{def:mpdo_triple_fusion_comparison,
        def:mpdo_final_label_selector_words}
    Suppose the positive trace-power coefficients are independent of the
    positive chain length and common word selectors exist.  If
    $\varepsilon\ne\varepsilon'$, then the corner of the full comparison from
    the right final sector $\varepsilon'$ to the left final sector
    $\varepsilon$ vanishes:
    \[
        P^{\mathrm L}_{\varepsilon}
        C_{\alpha,\beta,\gamma}
        P^{\mathrm R}_{\varepsilon'}=0.
    \]
    The assertion is conditional on the displayed word-selector identities.
    It neither states that a diagonal corner is invertible nor identifies one
    with an $F$-matrix, and it does not assert the pentagon identity.
\end{theorem}
\begin{proof}\leanok
    \uses{thm:mpdo_triple_fusion_comparison_intertwines}
    Fix one left multiplicity index and one right multiplicity index.  Taking
    the corresponding bond-space corner of the full letterwise identity gives
    \[
        M_\varepsilon^{ij}X
          =X M_{\varepsilon'}^{ij}.
    \]
    Induction on the word $w$ gives
    $M_\varepsilon^w X = X M_{\varepsilon'}^w$.  Multiply these identities by
    $c_\varepsilon(w)$ and sum over all words of length $S$.  The two selector
    identities give $X=0$.  Since the multiplicity indices were arbitrary,
    the whole off-diagonal final-sector corner vanishes.
\end{proof}

\begin{definition}[Fixed-final multiplicity spaces]%
    \label{def:mpdo_final_sector_multiplicity_spaces}
    \lean{MPOTensor.BNTFusionIsometryFamily.LeftFinalMultiplicity}
    \lean{MPOTensor.BNTFusionIsometryFamily.RightFinalMultiplicity}
    \lean{MPOTensor.BNTFusionIsometryFamily.leftFinalIndexEquiv}
    \lean{MPOTensor.BNTFusionIsometryFamily.rightFinalIndexEquiv}
    \lean{MPOTensor.BNTFusionIsometryFamily.leftTripleFinalEquiv}
    \lean{MPOTensor.BNTFusionIsometryFamily.rightTripleFinalEquiv}
    \leanok
    \uses{def:mpdo_bnt_fusion_isometry_family}
    Fix a final label $\varepsilon$.  The multiplicity spaces of the two
    bracketings are
    \[
        \mathcal H^{\mathrm L}_{\varepsilon}
          = \bigoplus_{\delta}
            \C^{\dim\chi_{\alpha,\beta,\delta}}
            \otimes \C^{\dim\chi_{\delta,\gamma,\varepsilon}},
        \qquad
        \mathcal H^{\mathrm R}_{\varepsilon}
          = \bigoplus_{\delta}
            \C^{\dim\chi_{\beta,\gamma,\delta}}
            \otimes \C^{\dim\chi_{\alpha,\delta,\varepsilon}}.
    \]
    These spaces have the bracketing pattern of the $(e,\mu,\nu)$ and
    $(f,\lambda,\sigma)$ index spaces in the printed-$F$ row/column convention
    of \cite[Section~3.4]{Bultinck2017Anyons}.  Identifying the dimensions of the positive
    diagonal matrices with the fusion multiplicities of that equation
    requires the length-independent integer specialization and is not
    asserted here.  The corresponding fixed-final row spaces are
    $\mathcal H^{\mathrm L}_{\varepsilon}\otimes\C^{D_\varepsilon}$ and
    $\mathcal H^{\mathrm R}_{\varepsilon}\otimes\C^{D_\varepsilon}$.
\end{definition}

\begin{theorem}[Equality of fixed-final multiplicity dimensions]%
    \label{thm:mpdo_final_sector_multiplicity_card_eq}
    \lean{MPOTensor.BNTFusionIsometryFamily.%
        card_leftFinalMultiplicity_eq_card_rightFinalMultiplicity_of_lengthIndependent}
    \leanok
    \uses{def:mpdo_final_sector_multiplicity_spaces,
        def:mpdo_bnt_label_linear_independent,
        def:mpdo_bnt_label_length_independent}
    Suppose that the positive trace-power coefficients are independent of
    the positive chain length and that the labelled operators are eventually
    linearly independent.  Then, for each final label $\varepsilon$,
    \[
        \dim\mathcal H^{\mathrm L}_{\varepsilon}
          = \sum_{\delta}
              \dim\chi_{\alpha,\beta,\delta}\,
              \dim\chi_{\delta,\gamma,\varepsilon}
          = \sum_{\delta}
              \dim\chi_{\beta,\gamma,\delta}\,
              \dim\chi_{\alpha,\delta,\varepsilon}
          = \dim\mathcal H^{\mathrm R}_{\varepsilon}.
    \]
    Consequently a matrix between the two fixed-final multiplicity spaces
    is square.  This is the numerical associativity statement of
    \cite[lines~237--241]{Bultinck2017Anyons}; it does not assert completeness
    of the fusion maps or invertibility of the fixed-final comparison.
\end{theorem}
\begin{proof}\leanok
    \uses{thm:mpdo_bnt_label_coeff_associativity,
        thm:mpdo_bnt_label_length_independent_natural}
    Choose a positive length beyond the linear-independence threshold.
    Associativity of the same-length product law gives equality of the two
    sums of products of structure coefficients.  Length independence and
    positivity identify every coefficient with the size of its corresponding
    diagonal $\chi$-matrix.  The two sums are precisely the cardinalities of
    the displayed dependent sums.
\end{proof}

\begin{definition}[Triple-fusion maps with fixed final label]%
    \label{def:mpdo_final_sector_fusion_maps}
    \lean{MPOTensor.BNTFusionIsometryFamily.leftFinalFusionMap}
    \lean{MPOTensor.BNTFusionIsometryFamily.rightFinalFusionMap}
    \leanok
    \uses{def:mpdo_final_sector_multiplicity_spaces,
        def:mpdo_left_fusion_isometry,
        def:mpdo_right_fusion_isometry}
    The maps $U^{\mathrm L}_{\alpha,\beta,\gamma;\varepsilon}$ and
    $U^{\mathrm R}_{\alpha,\beta,\gamma;\varepsilon}$ are obtained by
    restricting the two iterated fusion isometries to the summands with final
    label $\varepsilon$.  Their adjoints, after the canonical reassociation
    of the triple bond space, have the orientation and bracketing pattern of
    the weighted analogues of the two sides in the printed-$F$ row/column
    convention of \cite[Section~3.4]{Bultinck2017Anyons}.
\end{definition}

\begin{theorem}[Left triple fusion at a fixed final label]%
    \label{thm:mpdo_left_final_fusion_apply}
    \lean{MPOTensor.BNTFusionIsometryFamily.leftFinalFusion_apply}
    \leanok
    \uses{def:mpdo_final_sector_fusion_maps,
        def:mpdo_operator_product, thm:mpdo_left_fusion_apply,
        lem:matrix_block_conjugation_restriction}
    For every final label $\varepsilon$, the left-associated restriction
    satisfies the positive-diagonal weighted identity
    \[
        U^{\mathrm L}_{\alpha,\beta,\gamma;\varepsilon}
        \bigl((M_\alpha M_\beta)M_\gamma\bigr)^{ij}
        (U^{\mathrm L}_{\alpha,\beta,\gamma;\varepsilon})^\dagger
        = \bigoplus_{\delta}
          \chi_{\alpha,\beta,\delta}\otimes
          \bigl(\chi_{\delta,\gamma,\varepsilon}\otimes
          M_\varepsilon^{ij}\bigr).
    \]
\end{theorem}
\begin{proof}\leanok
    \uses{lem:matrix_block_conjugation_restriction}
    Restrict Theorem~\ref{thm:mpdo_left_fusion_apply} to the rows and columns
    whose final label is $\varepsilon$. Reindexing the conjugating matrix
    commutes with conjugation, and the outer block-diagonal matrix retains
    precisely the displayed $\varepsilon$-blocks.
\end{proof}

\begin{theorem}[Right triple fusion at a fixed final label]%
    \label{thm:mpdo_right_final_fusion_apply}
    \lean{MPOTensor.BNTFusionIsometryFamily.rightFinalFusion_apply}
    \leanok
    \uses{def:mpdo_final_sector_fusion_maps,
        def:mpdo_operator_product, thm:mpdo_right_fusion_apply,
        lem:matrix_block_conjugation_restriction}
    For every final label $\varepsilon$, the right-associated restriction
    satisfies the positive-diagonal weighted identity
    \[
        U^{\mathrm R}_{\alpha,\beta,\gamma;\varepsilon}
        \bigl(M_\alpha(M_\beta M_\gamma)\bigr)^{ij}
        (U^{\mathrm R}_{\alpha,\beta,\gamma;\varepsilon})^\dagger
        = \bigoplus_{\delta}
          \chi_{\beta,\gamma,\delta}\otimes
          \bigl(\chi_{\alpha,\delta,\varepsilon}\otimes
          M_\varepsilon^{ij}\bigr).
    \]
\end{theorem}
\begin{proof}\leanok
    \uses{lem:matrix_block_conjugation_restriction}
    Restrict Theorem~\ref{thm:mpdo_right_fusion_apply} to the rows and columns
    whose final label is $\varepsilon$. Reindexing the conjugating matrix
    commutes with conjugation, and the outer block-diagonal matrix retains
    precisely the displayed $\varepsilon$-blocks.
\end{proof}

\begin{theorem}[Injective fixed-final intertwiners]%
    \label{thm:mpdo_injective_fixed_final_intertwiner}
    \lean{MPOTensor.BNTFusionIsometryFamily.%
        exists_reindexed_intertwiner_eq_kronecker_one_of_isInjective}
    \leanok
    \uses{def:mpdo_final_sector_multiplicity_spaces, def:injective}
    Suppose that $M_\varepsilon$ is injective.  Let $C$ be a matrix from the
    right fixed-final-sector space to the left fixed-final-sector space.  Under
    the canonical identifications with
    $\mathcal H^{\mathrm L}_{\varepsilon}\otimes\C^{D_\varepsilon}$ and
    $\mathcal H^{\mathrm R}_{\varepsilon}\otimes\C^{D_\varepsilon}$, assume
    that, for every pair of physical indices $i,j$,
    \[
        (1_{\mathcal H^{\mathrm L}_{\varepsilon}}\otimes M_\varepsilon^{ij})C
        =C(1_{\mathcal H^{\mathrm R}_{\varepsilon}}\otimes
        M_\varepsilon^{ij}).
    \]
    Then there is a rectangular matrix
    $F:\mathcal H^{\mathrm R}_{\varepsilon}\to
    \mathcal H^{\mathrm L}_{\varepsilon}$ such that
    \[
        C=F\otimes 1_{D_\varepsilon}.
    \]
    This conclusion is conditional on the supplied matrix $C$.  It does not
    assert the existence or invertibility of $C$ or $F$, and it does not
    identify the dimensions of the positive diagonal matrices $\chi$ with the
    fusion multiplicities in the printed-$F$ row/column convention of
    \cite[Section~3.4]{Bultinck2017Anyons}.
\end{theorem}
\begin{proof}\leanok
    \uses{lem:amplified_center_scalar}
    Apply Lemma~\ref{lem:amplified_center_scalar} to the family
    $\{M_\varepsilon^{ij}\}_{i,j}$.  This family spans the full matrix algebra
    by injectivity, and the assumed relation is precisely the amplified
    intertwining identity of that lemma.
\end{proof}

\begin{theorem}[Conditional fixed-final Kronecker form and separation]%
    \label{thm:mpdo_triple_fusion_comparison_final_sector_kronecker}
    \lean{MPOTensor.BNTFusionIsometryFamily.%
        exists_tripleFusionComparison_finalSector_eq_kronecker_one_of_separation}
    \leanok
    \uses{def:mpdo_triple_fusion_comparison,
        def:mpdo_final_label_selector_words,
        def:mpdo_final_sector_multiplicity_spaces}
    Suppose the positive trace-power coefficients are independent of the
    positive chain length and common word selectors separate the final labels.
    Fix labels $\alpha,\beta,\gamma,\varepsilon$, and suppose that
    $M_\varepsilon$ is injective at the present blocking.  Write
    \[
        C^{\varepsilon,\varepsilon'}_{\alpha,\beta,\gamma}
          =P^{\mathrm L}_{\varepsilon}
            C_{\alpha,\beta,\gamma}
            P^{\mathrm R}_{\varepsilon'}
    \]
    for the corner from the right final sector $\varepsilon'$ to the left
    final sector $\varepsilon$.  There is a rectangular matrix
    \[
        F_\varepsilon:\mathcal H^{\mathrm R}_{\varepsilon}
          \longrightarrow \mathcal H^{\mathrm L}_{\varepsilon}
    \]
    such that, after the canonical index identifications,
    \[
        C^{\varepsilon,\varepsilon}_{\alpha,\beta,\gamma}
          =F_\varepsilon\otimes 1_{D_\varepsilon}.
    \]
    For every $\varepsilon'\ne\varepsilon$,
    \[
        C^{\varepsilon,\varepsilon'}_{\alpha,\beta,\gamma}=0.
    \]
    The selector identities and injectivity at the present blocking are
    additional hypotheses here.  The theorem neither asserts that
    $F_\varepsilon$ is square or invertible nor identifies it with the
    printed $F$-matrix of \cite[Section~3.4]{Bultinck2017Anyons}; no pentagon identity
    is asserted.
\end{theorem}
\begin{proof}\leanok
    \uses{thm:mpdo_triple_fusion_comparison_intertwines,
        thm:mpdo_triple_fusion_comparison_final_sector_vanishing,
        thm:mpdo_injective_fixed_final_intertwiner}
    Taking the $\varepsilon$--$\varepsilon$ corner of the full comparison
    identity gives, for every pair $i,j$,
    \[
        (1_{\mathcal H^{\mathrm L}_\varepsilon}\otimes
          M_\varepsilon^{ij})C^{\varepsilon,\varepsilon}
          =C^{\varepsilon,\varepsilon}
            (1_{\mathcal H^{\mathrm R}_\varepsilon}\otimes
              M_\varepsilon^{ij}).
    \]
    Injectivity and the amplified commutant result give
    $C^{\varepsilon,\varepsilon}=F_\varepsilon\otimes1_{D_\varepsilon}$.
    The common word selectors give
    $C^{\varepsilon,\varepsilon'}=0$ whenever
    $\varepsilon'\ne\varepsilon$ by
    Theorem~\ref{thm:mpdo_triple_fusion_comparison_final_sector_vanishing}.
\end{proof}

\begin{theorem}[Unitarity of each fixed-final comparison corner]%
    \label{thm:mpdo_triple_fusion_comparison_final_sector_unitary}
    \lean{MPOTensor.BNTFusionIsometryFamily.%
        tripleFusionComparison_finalSector_mul_conjTranspose_eq_one}
    \lean{MPOTensor.BNTFusionIsometryFamily.%
        conjTranspose_mul_tripleFusionComparison_finalSector_eq_one}
    \leanok
    \uses{def:mpdo_final_label_selector_words,
        def:mpdo_final_sector_multiplicity_spaces,
        def:mpdo_triple_fusion_comparison}
    Suppose the positive trace-power coefficients are independent of the
    positive chain length and common selectors exist at a positive length.
    For every final label $\varepsilon$, write

    \[
        C_\varepsilon
          =P^{\mathrm L}_{\varepsilon}
            C_{\alpha,\beta,\gamma}
            P^{\mathrm R}_{\varepsilon}.
    \]
    Then
    \begin{align}
        C_\varepsilon C_\varepsilon^\dagger&=1,
        \label{eq:mpdo_final_sector_right_unitary}\\
        C_\varepsilon^\dagger C_\varepsilon&=1.
        \label{eq:mpdo_final_sector_left_unitary}
    \end{align}
    The conclusion is conditional on the simultaneous selector identities.
    It is the fixed-output-channel invertibility consequence of the
    simultaneous inverse used in \cite[lines~247--277]{Bultinck2017Anyons}.
\end{theorem}
\begin{proof}\leanok
    \uses{thm:mpdo_triple_fusion_comparison_right_inverse,
        thm:mpdo_triple_fusion_comparison_left_inverse,
        thm:mpdo_triple_fusion_comparison_final_sector_vanishing}
    Order each full direct sum first by its final label.  In the product
    $CC^\dagger$, the $\varepsilon$-diagonal corner is
    \[
        \sum_{\varepsilon'}
          C^{\varepsilon,\varepsilon'}
          (C^{\varepsilon,\varepsilon'})^\dagger.
    \]
    Every term with $\varepsilon'\ne\varepsilon$ vanishes, while
    $CC^\dagger=1$.  Hence
    $C_\varepsilon C_\varepsilon^\dagger=1$.  Taking instead the
    $\varepsilon$-diagonal corner of $C^\dagger C=1$ gives the second identity.
\end{proof}
\begin{theorem}[Unitarity of the fixed-final multiplicity matrix]%
    \label{thm:mpdo_triple_fusion_multiplicity_unitary}
    \lean{MPOTensor.BNTFusionIsometryFamily.%
        exists_tripleFusionComparison_finalSector_eq_kronecker_one_and_unitary}
    \leanok
    \uses{def:mpdo_final_label_selector_words,
        def:mpdo_final_sector_multiplicity_spaces,
        def:mpdo_triple_fusion_comparison,
        lem:matrix_kronecker_identity_cancellation}
    Assume the hypotheses of
    Theorem~\ref{thm:mpdo_triple_fusion_comparison_final_sector_unitary}.
    Suppose in addition that $M_\varepsilon$ is injective at the present
    blocking and that $D_\varepsilon>0$.  There is a matrix
    \[
        F_\varepsilon:\mathcal H^{\mathrm R}_{\varepsilon}
          \longrightarrow \mathcal H^{\mathrm L}_{\varepsilon}
    \]
    such that
    \[
        C_\varepsilon=F_\varepsilon\otimes1_{D_\varepsilon},
        \qquad
        F_\varepsilon F_\varepsilon^\dagger=1,
        \qquad
        F_\varepsilon^\dagger F_\varepsilon=1.
    \]
    Thus $F_\varepsilon$ is invertible.  Its orientation is opposite to the
    printed-$F$ row/column convention of
    \cite[Section~3.4]{Bultinck2017Anyons}; no pentagon
    identity is asserted.
\end{theorem}
\begin{proof}\leanok
    \uses{thm:mpdo_triple_fusion_comparison_final_sector_kronecker,
        thm:mpdo_triple_fusion_comparison_final_sector_unitary,
        lem:matrix_kronecker_identity_cancellation}
    Injectivity gives
    $C_\varepsilon=F_\varepsilon\otimes1_{D_\varepsilon}$.
    Substitute this expression into equations~(\ref{eq:mpdo_final_sector_right_unitary})
    and (\ref{eq:mpdo_final_sector_left_unitary}):
    \[
        (F_\varepsilon F_\varepsilon^\dagger)
          \otimes1_{D_\varepsilon}=1,
        \qquad
        (F_\varepsilon^\dagger F_\varepsilon)
          \otimes1_{D_\varepsilon}=1.
    \]
    Since $D_\varepsilon>0$, Lemma~\ref{lem:matrix_kronecker_identity_cancellation}
    removes the identity factor and gives the two asserted identities.
\end{proof}

\begin{definition}[Simultaneous MPO block left inverse]
    \label{def:mpdo_block_left_inverse}
    \lean{MPSTensor.IsMPOBlockLeftInverse}
    \leanok
    Let $M_c^{il}\in M_{D_c}(\mathbb C)$ be a finite labelled MPO tensor
    family.  A simultaneous block left inverse consists of coefficients
    $C_{(c,x,y),(i,l)}$ such that
    \[
        \sum_{i,l}C_{(c,x,y),(i,l)}(M_d^{il})_{x'y'}
        =\delta_{cd}\delta_{xx'}\delta_{yy'}.
    \]
    This is the common block inverse used in
    \cite[Appendix~C.2, lines~1666--1676]{Cirac2016MPDO_arXiv} and
    \cite[lines~269--277]{Bultinck2017Anyons}.
\end{definition}

\begin{theorem}[Simultaneous inverse from a one-letter product-algebra span]%
    \label{thm:mpdo_one_letter_simultaneous_inverse}
    \lean{MPSTensor.WordTupleSpanTop.%
        exists_one_letter_simultaneous_left_inverse}
    \lean{MPSTensor.WordTupleSpanTop.exists_mpo_block_left_inverse}
    \leanok
    \uses{def:mpdo_block_left_inverse,
      thm:cpsv_bnt_blocked_simultaneous_span}
    Let $A_c^q\in M_{D_c}(\mathbb C)$ be a finite labelled tensor family.  If
    the one-letter tuples
    \[
        q\longmapsto (A_c^q)_c
    \]
    span the full product algebra $\prod_c M_{D_c}(\mathbb C)$, then there
    are coefficients $C_{(c,x,y),q}$ such that
    \[
        \sum_q C_{(c,x,y),q}(A_d^q)_{x'y'}
          =\delta_{cd}\delta_{xx'}\delta_{yy'}.
    \]
    For an MPO tensor the letter is the ket--bra pair $q=(i,k)$, so the same
    coefficients form the common block-letter left inverse used in
    \cite[lines~269--277]{Bultinck2017Anyons}.  The full product-algebra span
    is obtained after the common blocking described in
    \cite[lines~427--431]{Bultinck2017Anyons}.
\end{theorem}

\begin{proof}\leanok
    Fix $(c,x,y)$ and consider the product-algebra element which is the matrix
    unit $E_{xy}$ in block $c$ and zero in every other block.  By the spanning
    hypothesis it is a linear combination of the one-letter tuples.  Reading
    its $(x',y')$ entry in block $d$ gives the displayed identity.  Reindexing
    the doubled physical letter by the ket--bra pair gives the MPO form.
\end{proof}

\begin{figure}[ht]
    \centering
    % The first row is the word-selector identity itself: M^w is already the
    % matrix obtained by evaluating the entire word, so no local tensor
    % network is drawn.  In the second row the root labels epsilon and
    % epsilon' give the final-sector cut, and C and F map right-tree
    % coordinates to left-tree coordinates.  The last row is the
    % corresponding diagonal corner.
    \[
      \begin{aligned}
        \textup{selector},\ S>0:\quad &
        \sum_{|w|=S}c_{\varepsilon}(w)M_{\varepsilon'}^{w}
        =\delta_{\varepsilon,\varepsilon'}1_{D_\varepsilon}
        \\[1.5em]
        \textup{sector cut}:\quad &
        \underbrace{
          \tntree{((\alpha\,\beta)\,\gamma)_\varepsilon}
        }_{U^{\mathrm L}_{\alpha,\beta,\gamma;\varepsilon}}
        \xleftarrow{
          C^{\varepsilon,\varepsilon'}_{\alpha,\beta,\gamma}}
        \underbrace{
          \tntree{(\alpha\,(\beta\,\gamma))_{\varepsilon'}}
        }_{U^{\mathrm R}_{\alpha,\beta,\gamma;\varepsilon'}}
        =0\qquad(\varepsilon\ne\varepsilon')
        \\[1.5em]
        \textup{diagonal corner}:\quad &
        \mathcal H^{\mathrm L}_{\varepsilon}
          \otimes\C^{D_\varepsilon}
        \xleftarrow{F_\varepsilon\otimes1_{D_\varepsilon}}
        \mathcal H^{\mathrm R}_{\varepsilon}
          \otimes\C^{D_\varepsilon}
        \quad\Longrightarrow\quad
        F_\varepsilon F_\varepsilon^\dagger=1,\qquad
        F_\varepsilon^\dagger F_\varepsilon=1 .
      \end{aligned}
    \]
    \caption{The fixed-final comparison in the graphical order of
        \cite[equation~(pentagon3) and lines~269--277]
        {Bultinck2017Anyons}.  At a positive common length, the simultaneous
        word selector is the finite-word counterpart of the inverse
        $B_d^+$ in the source: it retains the chosen final label and
        annihilates every other one.  Thus the comparison has no
        off-diagonal final-label corners.  On the surviving corner,
        injectivity gives
        $C_\varepsilon=F_\varepsilon\otimes1_{D_\varepsilon}$; two-sided
        unitarity of $C_\varepsilon$ and $D_\varepsilon>0$ give two-sided
        unitarity of $F_\varepsilon$.  This is the associativity mechanism
        described in \cite[lines~995--1010]{Cirac2016MPDO_arXiv}.  No
        pentagon coherence identity is asserted here.}
    \label{fig:mpdo_fixed_final_comparison_unitary}
\end{figure}
